\documentclass[11pt,a4paper]{article}
\usepackage[T1]{fontenc}
\usepackage{lmodern}
\usepackage[margin=26mm]{geometry}
\usepackage{amsmath,amssymb,amsthm,booktabs,array,longtable}
\usepackage{microtype,multicol}
\usepackage[hidelinks]{hyperref}
\setlength{\parskip}{3pt}
\setlength{\emergencystretch}{2em}
\raggedbottom
\allowdisplaybreaks[1]
\clubpenalty=10000
\widowpenalty=10000
\newtheorem{theorem}{Theorem}[section]
\newtheorem{lemma}[theorem]{Lemma}
\newtheorem{corollary}[theorem]{Corollary}
\theoremstyle{remark}
\newtheorem{remark}[theorem]{Remark}
\newcommand{\Z}{\mathbb Z}
\newcommand{\Q}{\mathbb Q}
\newcommand{\N}{\mathbb N}
\newcommand{\F}{\mathbb F}
\newcommand{\ps}[1]{#1[[x]]}
\newcommand{\Rev}{\operatorname{Rev}}
\newcommand{\rev}{^{\langle-1\rangle}}
\newcommand{\Plus}{\operatorname{Plus}}
\newcommand{\Minus}{\operatorname{Minus}}
\newcommand{\coeff}[2]{[x^{#1}]#2}
\newcommand{\E}{\mathcal E}
\Urlmuskip=0mu plus 1mu
\hypersetup{pdftitle={A389536: Sparse residues in a Mahler substitution equation},pdfauthor={Dev Goyal},pdfsubject={Separate manuscript with AI disclosure and Lean proof map}}
\title{A389536\\[4pt]\large Sparse residues in a Mahler substitution equation}
\author{Dev Goyal}
\date{24 October 2026}
\begin{document}
\maketitle
\begin{center}
\fbox{\begin{minipage}{0.9\linewidth}\small
\textbf{AI disclosure.} ChatGPT (OpenAI) generated the mathematical arguments,
Lean proofs, and this manuscript in response to Dev Goyal's requests.
Dev Goyal is the sole human author. The source sequence and conjecture or
observation are credited to Paul D. Hanna.
\textbf{Independent human review is pending.}
\end{minipage}}
\end{center}
\begin{abstract}
The normalized solution of $A(x^2+2x^3)=xA(x)$ satisfies $A\equiv x+2xC\pmod4$, where $C=x+C^2$. Its coefficients above degree one are two modulo four exactly at indices $2^k+1$, and zero otherwise. The source equation is solved integrally before reduction.
\end{abstract}

\section{Source, notation, and scope}

The source is Paul D. Hanna's OEIS entry A389536 \cite{oeis536},
retrieved on 24 October 2026. Its initial sequence index is 1.
We also retain the formal zero constant coefficient at index zero.

All series are formal. The set $\N$ includes zero, and $[x^n]U$ denotes a
coefficient. A normalized series is $x+O(x^2)$; $O(x^r)$ means divisibility
by $x^r$, over any coefficient ring used below. We write $U\circ V=U(V(x))$
only when $V(0)=0$. Compositional inversion, denoted $\Rev(U)$ or $U\rev$,
is distinct from multiplicative inversion of a unit. Congruences are
coefficientwise; ordinary products and compositional iterates are distinct.

\section{Formal algebra used in the proof}

\begin{lemma}[Coefficientwise construction]\label{lem:contraction}
Let $K$ be a commutative ring. Suppose an operator $\Phi$ preserves
normalized series and, for $r\ge2$, satisfies
\[
 G-H=O(x^r)\quad\Longrightarrow\quad
 \Phi(G)-\Phi(H)=O(x^{r+1}).
\]
Then $\Phi$ has exactly one normalized fixed point in $\ps K$.
More generally, a residual equation whose degree-$r$ condition determines
the next unknown coefficient with multiplier one has a unique solution
once its initial coefficients are fixed and its lower-degree conditions
hold identically.
\end{lemma}

\begin{proof}
Begin with $G_0=x$ and set $G_{m+1}=\Phi(G_m)$. Normalization gives
$G_1-G_0=O(x^2)$. Inductively,
$G_{m+1}-G_m=O(x^{m+2})$, so each coefficient eventually stabilizes.
Define $G$ using those stable coefficients. The displayed implication
ensures that each coefficient of $\Phi(G)$ equals the corresponding stable
coefficient, hence $G=\Phi(G)$. If two fixed points agree below degree
$r$, they agree through degree $r$; induction proves uniqueness.

For the residual version, set the next coefficient to the negative of the
known residual term. Its multiplier is one, so the construction uses only
ring operations. No earlier equation changes, and every equation is
eventually imposed. The same coefficient comparison proves uniqueness.
\end{proof}

\begin{lemma}[Binary support]\label{lem:binary}
There is a unique normalized integral series $C$ satisfying $C=x+C^2$.
Its reduction modulo two is
\[
 P(x)=\sum_{k\ge0}x^{2^k}.
\]
More generally, for every integer $d\ge1$, the unique zero-constant
solution of $V+V^2=x^d$ in $\ps{\F_2}$ is
\[
 V=P(x^d)=\sum_{k\ge0}x^{d2^k}.
\]
\end{lemma}

\begin{proof}
The operator $G\mapsto x+G^2$ satisfies Lemma~\ref{lem:contraction},
since $G^2-H^2=(G-H)(G+H)$ and $G+H=O(x)$. In characteristic two,
every series $V$ satisfies $V(x)^2=V(x^2)$. This follows by truncating
above any fixed degree and applying the polynomial Frobenius identity.

Thus $V+V^2=x^d$ is a recursion in which the coefficient at an odd positive
index is fixed immediately, while the coefficient at an even index $n$
depends on the coefficient at $n/2<n$. The constant coefficient is fixed
to zero. This proves uniqueness. The proposed sum is well-defined
coefficientwise because its exponents tend to infinity; adding its square
cancels all terms except $x^d$. Set $d=1$ to identify the reduction of $C$.
\end{proof}

\begin{lemma}[Doubling and substitution]\label{lem:double}
For $U,V\in(\Z/4\Z)[[x]]$,
\[
 2U=2V\quad\Longleftrightarrow\quad
 U\equiv V\pmod2.
\]
If $B,D$ have zero constant coefficient and $B\equiv D\pmod2$, then
$2H(B)=2H(D)$ for every series $H$ over $\Z/4\Z$.
\end{lemma}

\begin{proof}
For residues $u,v$ modulo four, $2u=2v$ exactly when their reductions modulo
two agree. Apply this fact to each coefficient. Reduction modulo two
commutes with substitution, giving the second assertion. In particular,
this lemma never cancels two inside $\Z/4\Z$.
\end{proof}

\section{A389536: a Mahler-type substitution modulo four}

\begin{theorem}\label{thm:mahler}
The source equation
\begin{equation}\label{eq:mahler}
 A(x^2+2x^3)=xA(x),\qquad A=x+O(x^2),
\end{equation}
has a unique solution in $\ps{\Z}$. With $C=x+C^2$ as in
Lemma~\ref{lem:binary}, it satisfies
\[
 A(x)\equiv x+2xC(x)\pmod4.
\]
Consequently, for every $n>1$,
\[
 [x^n]A\equiv
 \begin{cases}
  2\pmod4,&n=2^k+1\text{ for some }k\in\N,\\
  0\pmod4,&\text{otherwise}.
 \end{cases}
\]
\end{theorem}

\begin{proof}
Write $B=x^2+2x^3$. For $A=\sum_{n\ge1}a_nx^n$ with $a_1=1$, the
coefficient of $x^{n+1}$ in \eqref{eq:mahler}, for $n\ge2$, gives
\begin{equation}\label{eq:mahlerrec}
 a_n=
 \sum_{\substack{j\ge1\\0\le n+1-2j\le j}}
 a_j\binom{j}{\,n+1-2j\,}2^{n+1-2j}.
\end{equation}
Indeed, $B^j=x^{2j}(1+2x)^j$. Every contributing index obeys
$j\le(n+1)/2<n$, so this is a finite, division-free recursion in earlier
coefficients. The degree-two equation is already satisfied by $a_1=1$,
and the lower equations vanish. Thus the construction exists and is
unique over any commutative ring and commutes with reduction.

Work now modulo four. Denote the reduction of $C$ by the same letter,
so $C=x+C^2$, and put $D=x+2xC$. Its normalization follows from $C(0)=0$.
Modulo two, $B=x^2$, and Frobenius gives $C(B)=C(x^2)=C^2$.
Lemma~\ref{lem:double} therefore gives $2C(B)=2C^2$ modulo four. Compute
in $(\Z/4\Z)[[x]]$:
\begin{align*}
 D(B)
 &=B+2B\,C(B)\\
 &=x^2+2x^3+2x^2C(B)\\
 &=x^2+2x^3+2x^2C^2\\
 &=x^2+2x^2C=xD.
\end{align*}
The second line discards a multiple of four; the fourth uses $C^2=C-x$.
Thus $D$ satisfies the actual source equation in this ring. Uniqueness
identifies it with the reduction of the integer solution.

For $n>1$, its coefficient is $2[x^{n-1}]C$ modulo four.
Lemma~\ref{lem:binary} says this is two precisely when $n-1=2^k$;
otherwise it is zero.
\end{proof}

The source writes $A=A(x^2+2x^3)/x$ \cite{oeis536}.
Equation~\eqref{eq:mahler} is the equivalent division-free identity.
Multiplication by $x$ is injective in every formal power-series ring,
so this rewriting introduces no additional choice of solution.


\section{Separate Lean proof and checked scope}

The entry file is \nolinkurl{lean/Entries/A389536.lean}. Its explicit
theorem statements and proof terms use the namespace
\nolinkurl{FinalCollection.A389536}. The complete imported implementations
and their shared mathematical dependencies are included in the collection.
The following local theorem names give the proof map:

\begin{multicols}{2}\raggedright\begin{itemize}\setlength{\itemsep}{2pt}\setlength{\parsep}{0pt}
\item \nolinkurl{mahler_existence}.
\item \nolinkurl{mahler_definition}.
\item \nolinkurl{mahler_source_unique}.
\item \nolinkurl{mahler_mod_four}.
\item \nolinkurl{mahler_residue}.
\item \nolinkurl{mahler_four_dvd_iff}.
\end{itemize}\end{multicols}

The main implementation module is \nolinkurl{Round3.MahlerParity}.
The project pins Lean 4.19.0 and Mathlib
\nolinkurl{c44e0c8ee63ca166450922a373c7409c5d26b00b}.
The supplied collection report records fresh source compilation and
transitive axiom audits. All coefficient conclusions are unbounded.

\paragraph{Limits and status.}
The source quotient by x is expressed by the equivalent division-free equation. Analytic claims in the entry are outside scope.
No global novelty, exclusive priority, peer review, or OEIS acceptance is asserted.

\begin{thebibliography}{9}
\raggedright
\bibitem{oeis536}
Paul D. Hanna, \emph{OEIS A389536}, 28 October 2025.
\url{https://oeis.org/A389536}. Retrieved 24 October 2026.
\end{thebibliography}
\end{document}
